% !TeX spellcheck = en_GB
\chapter*{Introduction}
\addcontentsline{toc}{chapter}{Introduction}

Evolutionary algorithms (EAs) are a class of optimization and search techniques inspired by the principles of natural evolution. Over the past decades, these algorithms have been successfully applied to a wide range of complex problems in engineering, robotics, and machine learning. Their ability to explore large and poorly understood search spaces makes them particularly suitable for problems where traditional optimization methods are applied only with great difficulty.

The performance of evolutionary algorithms is typically driven by a objective function, which evaluates candidate solutions according to predefined objectives. In classical evolutionary computation, this is usually refered to as fitness-based selection. While this approach has proven effective in many domains, it also presents several limitations. In particular, fitness-driven search can suffer from premature convergence, stagnation in local optima, and deception, where the search process becomes trapped in regions that appear promising but ultimately prevent discovery of globally optimal solutions.

To address these challenges, alternative search paradigms have been proposed, among which novelty search has gained significant attention. Unlike traditional fitness-based optimization, novelty search rewards behavioral uniqueness rather than objective performance. Introduced as a method to overcome deception in evolutionary search, novelty search encourages exploration by promoting individuals whose behaviors differ from those previously encountered. This shift from objective-oriented optimization to exploration-oriented search has demonstrated promising results in domains characterized by sparse rewards, deceptive landscapes, and high-dimensional complexity.

The distinction between fitness-based and novelty-driven approaches reflects two fundamentally different philosophies of evolutionary search. Fitness-oriented methods prioritize exploitation of known promising regions in the search space, whereas novelty-based methods emphasize exploration and behavioral diversity. Both approaches offer unique advantages and disadvantages depending on the characteristics of the problem domain. Fitness-based methods generally provide faster convergence toward measurable objectives but may lack robustness in deceptive environments. Novelty search, on the other hand, can discover innovative and unexpected solutions, although it may require greater computational resources and may struggle in tasks where exploration alone is insufficient.

This thesis focuses on a comparative analysis of novelty and fitness mechanisms in evolutionary algorithms when applied to reinforcement learning tasks. The objective is to investigate how these approaches influence search dynamics´, convergence behavior, diversity maintenance, and solution quality across selected optimization problems. The study aims to identify conditions under which novelty search outperforms traditional fitness-based evolution, as well as scenarios where fitness-driven optimization remains more effective.

The research combines theoretical analysis with experimental evaluation. Several benchmark problems and simulation environments are employed to compare the behavior of both methods under controlled conditions. Metrics such as convergence speed, population diversity, robustness, and final solution quality are analyzed to provide a comprehensive understanding of the strengths and limitations of each approach. Additionally, hybrid methods that combine novelty and fitness objectives are considered to examine whether balancing exploration and exploitation can improve overall evolutionary performance.

The findings of this thesis contribute to the broader understanding of exploration and exploitation trade-offs in evolutionary computation. By examining the relationship between novelty and fitness, this work seeks to provide insights into the design of more adaptive and efficient evolutionary algorithms capable of solving increasingly complex optimization problems.
